\documentclass[11pt,fontset=fandol]{ctexart}
\usepackage[a4paper,margin=23mm]{geometry}
\usepackage{amsmath,amssymb,mathtools,bm,mathrsfs,esint}
\usepackage{graphicx,longtable,booktabs,array,calc}
\usepackage[normalem]{ulem}
\usepackage{xcolor,hyperref}
\hypersetup{hidelinks}
\setcounter{secnumdepth}{-1}
\newcounter{none}
\providecommand{\tightlist}{\setlength{\itemsep}{0pt}\setlength{\parskip}{0pt}}
\providecommand{\passthrough}[1]{#1}
\usepackage{amssymb,mathrsfs,esint,bm,newunicodechar}
\allowdisplaybreaks
\setlength{\emergencystretch}{3em}
\providecommand{\llangle}{\langle\!\langle}
\providecommand{\rrangle}{\rangle\!\rangle}
\newunicodechar{∞}{\ensuremath{\infty}}
\newunicodechar{∈}{\ensuremath{\in}}
\newunicodechar{∑}{\ensuremath{\sum}}
\newunicodechar{∫}{\ensuremath{\int}}
\newunicodechar{≤}{\ensuremath{\leq}}
\newunicodechar{≥}{\ensuremath{\geq}}
\newunicodechar{≠}{\ensuremath{\neq}}
\newunicodechar{→}{\ensuremath{\to}}
\newunicodechar{⇒}{\ensuremath{\Rightarrow}}
\newunicodechar{α}{\ensuremath{\alpha}}
\newunicodechar{β}{\ensuremath{\beta}}
\newunicodechar{γ}{\ensuremath{\gamma}}
\newunicodechar{λ}{\ensuremath{\lambda}}
\newunicodechar{φ}{\ensuremath{\varphi}}
\newunicodechar{ε}{\ensuremath{\varepsilon}}
\newunicodechar{σ}{\ensuremath{\sigma}}
\newunicodechar{θ}{\ensuremath{\theta}}
\newunicodechar{δ}{\ensuremath{\delta}}
\newunicodechar{η}{\ensuremath{\eta}}
\newunicodechar{π}{\ensuremath{\pi}}
\newunicodechar{ξ}{\ensuremath{\xi}}
\newunicodechar{Σ}{\ensuremath{\Sigma}}
\newunicodechar{Δ}{\ensuremath{\Delta}}
\newunicodechar{∂}{\ensuremath{\partial}}
\newunicodechar{∀}{\ensuremath{\forall}}
\newunicodechar{∃}{\ensuremath{\exists}}
\newunicodechar{√}{\ensuremath{\surd}}
\newunicodechar{−}{\ensuremath{-}}
\newunicodechar{μ}{\ensuremath{\mu}}
\newunicodechar{ρ}{\ensuremath{\rho}}
\newunicodechar{τ}{\ensuremath{\tau}}
\newunicodechar{ω}{\ensuremath{\omega}}
\newunicodechar{Ω}{\ensuremath{\Omega}}
\newunicodechar{χ}{\ensuremath{\chi}}
\newunicodechar{ψ}{\ensuremath{\psi}}
\newunicodechar{Γ}{\ensuremath{\Gamma}}
\newunicodechar{Φ}{\ensuremath{\Phi}}
\newunicodechar{Λ}{\ensuremath{\Lambda}}
\newunicodechar{ℝ}{\ensuremath{\mathbb{R}}}
\newunicodechar{ℤ}{\ensuremath{\mathbb{Z}}}
\newunicodechar{ℕ}{\ensuremath{\mathbb{N}}}
\newunicodechar{△}{\ensuremath{\triangle}}
\newunicodechar{⊗}{\ensuremath{\otimes}}
\begin{document}
\section{2023年高等代数真题}\label{ux5e74ux9ad8ux7b49ux4ee3ux6570ux771fux9898}

来源：\texttt{HDU/601数学分析/2023/2023年杭州电子科技大学数分高代真题及答案(双科合卷).pdf}。原
PDF 页序：1, 2。

\section{杭州电子科技大学2023年高等代数真题}\label{ux676dux5ddeux7535ux5b50ux79d1ux6280ux5927ux5b662023ux5e74ux9ad8ux7b49ux4ee3ux6570ux771fux9898}

\begin{quote}
整理说明：保留全部九题及解析的数学内容，清理公式中的水印并恢复矩阵行列。第三题的系数阵经同年手写卷本地OCR核对，主对角线全为
\(\lambda\)、其余元素全为
\(1\)。第八题（2）的外层求和下标按公式和解析中的逐项恒等式恢复为
\(i\)，固定整数参数仍为
\(k\)。手写异文另附，不将无法辨识的手写内容混入本打印题面。
\end{quote}

\subsection{一、计算行列式}\label{ux4e00ux8ba1ux7b97ux884cux5217ux5f0f}

\[
D=\begin{vmatrix}
1&2&3&4&5\\2&3&4&5&1\\3&4&5&1&2\\4&5&1&2&3\\5&1&2&3&4
\end{vmatrix}.
\]

\subsection{二、矩阵方程}\label{ux4e8cux77e9ux9635ux65b9ux7a0b}

设 \[
A=\begin{pmatrix}1&1&0&-3\\4&-2&6&3\end{pmatrix}.
\]

（1）求 \(4\times2\) 矩阵 \(B\)，使 \(AB=0\) 且 \(r(B)=2\)。

（2）求矩阵 \(X\)，使 \(AX=E\)，其中 \(E\) 为单位矩阵。

\subsection{三、线性方程组（20分）}\label{ux4e09ux7ebfux6027ux65b9ux7a0bux7ec420ux5206}

\[
\begin{cases}
\lambda x_1+x_2+x_3+x_4=1,\\
x_1+\lambda x_2+x_3+x_4=1,\\
x_1+x_2+\lambda x_3+x_4=\lambda,\\
x_1+x_2+x_3+\lambda x_4=\lambda.
\end{cases}
\]

（1）\(\lambda\) 为何值时方程组有唯一解？求其唯一解。

（2）\(\lambda\) 为何值时方程组有无穷多解？求其通解。

\subsection{四、二次型（15分）}\label{ux56dbux4e8cux6b21ux578b15ux5206}

设 \[
f(x_1,x_2,x_3)=2x_1^2+5x_2^2+5x_3^2+4x_1x_2-4x_1x_3-8x_2x_3.
\] 用正交变换 \(X=TY\) 将二次型化为标准形，并求正交矩阵 \(T\)。

\subsection{五、线性变换（20分）}\label{ux4e94ux7ebfux6027ux53d8ux636220ux5206}

设 \(\varepsilon_1,\varepsilon_2,\varepsilon_3\) 是三维线性空间 \(V\)
的一组基，线性变换 \(\mathcal A\) 在该基下的矩阵为 \[
A=\begin{pmatrix}1&0&2\\-1&2&1\\1&2&5\end{pmatrix}.
\]

（1）求 \(\mathcal A\) 在基 \[
\eta_1=\varepsilon_1-2\varepsilon_2,\quad
\eta_2=3\varepsilon_2-\varepsilon_3,\quad\eta_3=\varepsilon_3
\] 下的矩阵 \(B\)。

（2）求 \(\mathcal AV\) 与 \(\mathcal A^{-1}(\theta)\)，其中 \(\theta\)
为零向量。

\subsection{六、矩阵求逆（15分）}\label{ux516dux77e9ux9635ux6c42ux900615ux5206}

已知 \(A,B,C\) 均为 \(n\) 阶方阵，\(E\)
为单位矩阵，\(D=E+BCA\)，\(C=(E-AB)^{-1}\)。证明：

（1）\(E-BA\) 可逆，且 \(D=(E-BA)^{-1}\)。

（2）\(C=E+ADB\)。

\subsection{七、公共特征向量（15分）}\label{ux4e03ux516cux5171ux7279ux5f81ux5411ux91cf15ux5206}

设 \(A,B\) 为复 \(n\) 维线性空间 \(V\) 上的线性变换，且
\(AB-BA=A\)。证明：

（1）\(A\) 的所有特征值为零。

（2）\(A\) 与 \(B\) 必有公共特征向量。

\subsection{八、实根多项式（15分）}\label{ux516bux5b9eux6839ux591aux9879ux5f0f15ux5206}

设实系数多项式 \(f(x)=\prod_{i=1}^n(x-a_i)\)，其中 \(a_1,\ldots,a_n\)
是两两互异的实数。证明：

（1）多项式 \([f'(x)]^2-f(x)f''(x)\) 没有实根。

（2） \[
f'(x)=\sum_{i=1}^n\left[\frac1{x-a_i}\left(\frac{a_i}{x}\right)^{k+1}
+\frac1x\sum_{r=0}^k\left(\frac{a_i}{x}\right)^r\right]f(x),\qquad k\in\mathbb Z,\quad k\ge2.
\] 该有分母的表达式先在 \(x\ne0,a_1,\ldots,a_n\)
处理解；约去可去因子后的多项式恒等式可延续到其他点。

\subsection{九、根子空间分解（20分）}\label{ux4e5dux6839ux5b50ux7a7aux95f4ux5206ux89e320ux5206}

设 \(A\) 是数域 \(P\) 上 \(n\) 维线性空间 \(V\)
上的线性变换，其特征多项式为 \[
f(\lambda)=(\lambda-\lambda_1)^{r_1}(\lambda-\lambda_2)^{r_2}\cdots(\lambda-\lambda_s)^{r_s}.
\] 其中 \(\lambda_i\in P\)
两两互异（同年手写卷末尾也保留了``两两互异''）。证明 \[
V=V_1\oplus\cdots\oplus V_s,\qquad
V_i=\{\xi\in V:(\lambda_iE-A)^{r_i}\xi=\theta\},\quad i=1,\ldots,s,
\] \(E\) 为恒等变换，\(\theta\) 为零向量。
\end{document}
